\documentclass{beamer}
\usepackage[utf8]{inputenc}

\usetheme{Madrid}
\usecolortheme{default}
\usepackage{amsmath,amssymb,amsfonts,amsthm}
\usepackage{txfonts}
\usepackage{tkz-euclide}
\usepackage{listings}
\usepackage{adjustbox}
\usepackage{array}
\usepackage{tabularx}
\usepackage{gvv}
\usepackage{lmodern}
\usepackage{circuitikz}
\usepackage{tikz}
\usepackage{graphicx}
\usepackage{multicol}

\setbeamertemplate{page number in head/foot}[totalframenumber]

\usepackage{tcolorbox}
\tcbuselibrary{minted,breakable,xparse,skins}



\definecolor{bg}{gray}{0.95}
\DeclareTCBListing{mintedbox}{O{}m!O{}}{%
  breakable=true,
  listing engine=minted,
  listing only,
  minted language=#2,
  minted style=default,
  minted options={%
    linenos,
    gobble=0,
    breaklines=true,
    breakafter=,,
    fontsize=\small,
    numbersep=8pt,
    #1},
  boxsep=0pt,
  left skip=0pt,
  right skip=0pt,
  left=25pt,
  right=0pt,
  top=3pt,
  bottom=3pt,
  arc=5pt,
  leftrule=0pt,
  rightrule=0pt,
  bottomrule=2pt,
  toprule=2pt,
  colback=bg,
  colframe=orange!70,
  enhanced,
  overlay={%
    \begin{tcbclipinterior}
    \fill[orange!20!white] (frame.south west) rectangle ([xshift=20pt]frame.north west);
    \end{tcbclipinterior}},
  #3,
}
\lstset{
    language=C,
    basicstyle=\ttfamily\small,
    keywordstyle=\color{blue},
    stringstyle=\color{orange},
    commentstyle=\color{green!60!black},
    numbers=left,
    numberstyle=\tiny\color{gray},
    breaklines=true,
    showstringspaces=false,
}
%------------------------------------------------------------
%This block of code defines the information to appear in the
%Title page
\title %optional
{5.8.27}
\date{October 12,2025}


\author 
{Jnanesh Sathisha karmar - EE25BTECH11029}



\begin{document}



\frame{\titlepage}
\begin{frame}{Question:}


\noindent $\vec{A}=\myvec{1 & 0 & 0\\0 & 1 & 1\\0 & -2 & 4}$ and $\vec{I}=\myvec{1 & 0 & 0\\0 & 1 & 0\\0 & 0 &1}$ and $\vec{A}^{-1} = \frac{1}{6}(\vec{A}^2 + c\vec{A} + d\vec{I})$, then value of c and d are
\begin{enumerate}
\begin{multicols}{4}
    \item \brak{-6,-11}
    \item \brak{6,11}
    \item \brak{-6,11}
    \item \brak{6,-11}
\end{multicols}
\end{enumerate}
\end{frame}

\begin{frame}{Theoretical Solution}
The characteristic equation for $\vec{A}$is given by $\det(A - \lambda I) = 0$.

First, we construct the matrix $\brak{\vec{A} - \lambda \vec{I}}$:
\begin{align}
\vec{A} - \lambda \vec{I} = \myvec{ 1 & 0 & 0 \\ 0 & 1 & 1 \\ 0 & -2 & 4 } - \lambda \myvec{ 1 & 0 & 0 \\ 0 & 1 & 0 \\ 0 & 0 & 1 } = \myvec{ 1-\lambda & 0 & 0 \\ 0 & 1-\lambda & 1 \\ 0 & -2 & 4-\lambda }
\end{align}
Next, we calculate its determinant:
\begin{align*}
\det(\vec{A} - \lambda \vec{I}) &= \brak{1-\lambda} \myvec{ 1-\lambda & 1 \\ -2 & 4-\lambda } \\
&= \brak{1-\lambda} \sbrak{(1-\lambda)(4-\lambda) - (1)(-2) } \\
&= \brak{1-\lambda} \sbrak{ 4 - 4\lambda - \lambda + \lambda^2 + 2 } \\
&= \brak{1-\lambda} \brak{ \lambda^2 - 5\lambda + 6 } \\
&= \lambda^2 - 5\lambda + 6 - \lambda^3 + 5\lambda^2 - 6\lambda \\
&= -\lambda^3 + 6\lambda^2 - 11\lambda + 6
\end{align*}
\end{frame}
\begin{frame}{Theoretical Solution}
Setting the determinant to zero gives the characteristic equation. Multiplying by -1 for convention, we get:
\begin{align}
\lambda^3 - 6\lambda^2 + 11\lambda - 6 = 0
\end{align}
The theorem allows us to substitute the matrix $A$ for $\lambda$ in the characteristic equation. The constant term is multiplied by the identity matrix $I$.
\begin{align}
\vec{A}^3 - 6\vec{A}^2 + 11\vec{A} - 6\vec{I} = \vec{O}
\end{align}
where $\vec{O}$ is the zero matrix.\\
We can rearrange Equation \brak{4} to solve for the inverse. First, move the identity term to the right side:
\begin{align}
\vec{A}^3 - 6\vec{A}^2 + 11\vec{A} = 6\vec{I}
\end{align}
Now, multiply the entire equation by $\vec{A}^{-1}$ from the left:
\begin{align*}
\vec{A}^{-1}\brak{\vec{A}^3 - 6\vec{A}^2 + 11\vec{A}} &= \vec{A}^{-1}\brak{6\vec{I}} \\
\vec{A}^2 - 6\vec{A} + 11\vec{I} &= 6\vec{A}^{-1}
\end{align*}
\end{frame}
\begin{frame}{Theoretical Solution}
Finally, isolate $\vec{A}^{-1}$ by dividing by 6:
\begin{align}
\vec{A}^{-1} = \frac{1}{6}\brak{\vec{A}^2 - 6\vec{A} + 11\vec{I}}
\end{align}
We now compare the derived expression \brak{5} with the expression given in the question.\\
Derived Expression: $\vec{A}^{-1} = \frac{1}{6}\brak{\vec{A}^2 - 6\vec{A} + 11\vec{I}}$\\
Given Expression: $\vec{A}^{-1} = \frac{1}{6}(\vec{A}^2 + c\vec{A} + d\vec{I})$\\
By directly comparing the coefficients of the matrices $\vec{A}$ and $\vec{I}$ inside the parentheses, we find:
\begin{align}
c = -6 \quad \text{and} \quad d = 11
\end{align}
The solution is therefore third option, $\brak{c, d} = \brak{-6, 11}$.
\end{frame}

\begin{frame}[fragile]
    \frametitle{C Code }

    \begin{lstlisting}
#include <stdio.h>
#include <stdlib.h>
#include <stdbool.h>
#include <math.h>

#define N 3

void printMatrix(double mat[N][N], const char* name);
void multiplyMatrices(double mat1[N][N], double mat2[N][N], double result[N][N]);
bool invertMatrix(double mat[N][N], double inverse[N][N]);
void scalarMultiply(double scalar, double mat[N][N], double result[N][N]);
void subtractMatrices(double mat1[N][N], double mat2[N][N], double result[N][N]);

int main() {
    double A[N][N] = {
        {1.0, 0.0, 0.0},
        {0.0, 1.0, 1.0},
        {0.0, -2.0, 4.0}
    };

    double A_squared[N][N];
    double A_inverse[N][N];
    double temp_matrix[N][N];
    double RHS_matrix[N][N];

    \end{lstlisting}
\end{frame}

\begin{frame}[fragile]
    \frametitle{C Code }
    \begin{lstlisting}
    printf("--- Initial Matrix A ---\n");
    printMatrix(A, "A");
    multiplyMatrices(A, A, A_squared);
    printf("\n--- Step 1: Calculated A^2 ---\n");
    printMatrix(A_squared, "A^2");
    printf("\n--- Step 2: Calculating A^-1 ---\n");
    if (!invertMatrix(A, A_inverse)) {
        printf("Matrix A is not invertible. Cannot solve.\n");
        return 1;}
    printMatrix(A_inverse, "A^-1");
    printf("\n--- Step 3: Calculating RHS = 6*A^-1 - A^2 ---\n");
    scalarMultiply(6.0, A_inverse, temp_matrix);
    printMatrix(temp_matrix, "6 * A^-1");
    subtractMatrices(temp_matrix, A_squared, RHS_matrix);
    printMatrix(RHS_matrix, "RHS Matrix");
    double c = RHS_matrix[1][2];
    double d = RHS_matrix[1][1] - c;
\end{lstlisting}
\end{frame}
\begin{frame}[fragile]
    \frametitle{C Code }
    \begin{lstlisting}
       printf("\n--- Step 4: Solved for c and d ---\n");
    printf("From element (2, 3) of the matrix equation, we get c = %.2f\n", c);
    printf("From element (2, 2), we get c + d = %.2f, so d = %.2f - %.2f = %.2f\n", RHS_matrix[1][1], RHS_matrix[1][1], c, d);
    printf("\n====================================\n");
    printf("           Final Solution           \n");
    printf("====================================\n");
    printf("c = %.2f\n", c);
    printf("d = %.2f\n", d);
    printf("====================================\n");
    return 0;}
void printMatrix(double mat[N][N], const char* name) {
    printf("Matrix %s:\n", name);
    for (int i = 0; i < N; i++) {
        for (int j = 0; j < N; j++) {printf("%8.3f ", mat[i][j]);}
        printf("\n");
    }
}
    
\end{lstlisting}
\end{frame}
\begin{frame}[fragile]
    \frametitle{C Code }
    \begin{lstlisting}
void multiplyMatrices(double mat1[N][N], double mat2[N][N], double result[N][N]) {
    for (int i = 0; i < N; i++) {
        for (int j = 0; j < N; j++) {
            result[i][j] = 0;
            for (int k = 0; k < N; k++) {
                result[i][j] += mat1[i][k] * mat2[k][j];
            }
        }
    }
}

bool invertMatrix(double mat[N][N], double inverse[N][N]) {
    double augvec[N][N * 2];

    for (int i = 0; i < N; i++) {
        for (int j = 0; j < N; j++) {
            augvec[i][j] = mat[i][j];
        }
        for (int j = 0; j < N; j++) {
            augvec[i][j + N] = (i == j) ? 1.0 : 0.0;
        }
    }
\end{lstlisting}
\end{frame}
\begin{frame}[fragile]
    \frametitle{C Code }
    \begin{lstlisting}
        for (int i = 0; i < N; i++) {
        double pivot = augvec[i][i];

        if (pivot == 0) {
            int max_row = i;
            for (int k = i + 1; k < N; k++) {
                if (fabs(augvec[k][i]) > fabs(augvec[max_row][i])) {
                    max_row = k;
                }
            }

            if (max_row != i) {
                for (int j = 0; j < 2 * N; j++) {
                    double temp = augvec[i][j];
                    augvec[i][j] = augvec[max_row][j];
                    augvec[max_row][j] = temp;
                }
            }
        }
\end{lstlisting}
\end{frame}
\begin{frame}[fragile]
    \frametitle{C Code }
    \begin{lstlisting}
            pivot = augvec[i][i];
        if (pivot == 0) return false;
        for (int j = 0; j < 2 * N; j++) {augvec[i][j] /= pivot;}
        for (int k = 0; k < N; k++) {
            if (k != i) {
                double factor = augvec[k][i];
                for (int j = 0; j < 2 * N; j++) {
                    augvec[k][j] -= factor * augvec[i][j];}
            }
        }
    }
    for (int i = 0; i < N; i++) {
        for (int j = 0; j < N; j++) {
            inverse[i][j] = augvec[i][j + N];
        }
    }
    return true;
}
\end{lstlisting}
\end{frame}
\begin{frame}[fragile]
    \frametitle{C Code }
    \begin{lstlisting}
    void scalarMultiply(double scalar, double mat[N][N], double result[N][N]) {
    for (int i = 0; i < N; i++) {
        for (int j = 0; j < N; j++) {
            result[i][j] = mat[i][j] * scalar;
        }
    }
}

void subtractMatrices(double mat1[N][N], double mat2[N][N], double result[N][N]) {
    for (int i = 0; i < N; i++) {
        for (int j = 0; j < N; j++) {
            result[i][j] = mat1[i][j] - mat2[i][j];
        }
    }
}
\end{lstlisting}
\end{frame}




\end{document}
